\documentclass[11pt,a4paper]{article}
\usepackage[utf8]{inputenc}
\usepackage[T1]{fontenc}
\usepackage{times}
\usepackage{amsmath,amssymb}
\usepackage{booktabs}
\usepackage{multirow}
\usepackage{graphicx}
\usepackage{xcolor}
\usepackage{hyperref}
\usepackage{natbib}
\usepackage[margin=2.5cm]{geometry}
\usepackage{caption}
\usepackage{subcaption}

\newcommand{\rmark}[1]{\textcolor{red}{#1}}
\newcommand{\bmark}[1]{\textbf{#1}}

\title{The Constrained Decoding Bottleneck: Why Confidence Signals\\Predict but Cannot Select in Information Extraction}

\author{Anonymous}

\date{}

\begin{document}
\maketitle

% ============================================================
\begin{abstract}
Test-time scaling via multi-sample generation improves LLM performance on reasoning tasks, but its effectiveness for structured information extraction (IE) remains unclear. We systematically evaluate five confidence signals---spanning log-probability, exact match, set similarity, token voting, and n-gram overlap---for selecting the best output among multiple candidates generated by fine-tuned LLMs on named entity recognition (NER) and relation extraction (RE). We identify a \emph{regime-specific selection bottleneck}: while these signals reliably estimate output quality (AUROC up to 0.75), they fail to select better candidates than greedy decoding under grammar-constrained generation ($\Delta$F1 $= +0.6$pp, $p = 0.39$). Diagnosis reveals that constrained decoding compresses log-probability variance---60\% of instances have within-sample log-prob range $< 0.05$---destroying the only per-sample ranking signal. A controlled ablation confirms the causal mechanism: removing grammar constraints recovers log-prob variance by 10$\times$ and enables significant selection gains (NER $+5.0$pp, $p = 0.008$). The signal ranking itself reverses: log-prob shifts from worst under constrained decoding ($\rho = 0.21$) to best under free-form generation ($\rho = 0.42$). A supervised DeBERTa verifier achieves only $+3.6$pp under constrained decoding, below the zero-cost free-form log-prob gain, confirming that the bottleneck lies in the decoding regime, not the scoring function. Our results provide the first controlled evidence that constrained decoding creates a fundamental selection barrier for structured IE, and suggest relaxed decoding with post-hoc validation as a practical mitigation path.
\end{abstract}

% ============================================================
\section{Introduction}
\label{sec:intro}

Multi-sample generation at test time has emerged as a powerful strategy for improving LLM outputs. Self-consistency \citep{wang2023selfconsistency} selects answers by majority vote, MBR decoding \citep{eikema2022sampling} selects translations by consensus similarity, and best-of-N sampling with reward models \citep{cobbe2021training,lightman2023lets} selects reasoning chains by learned quality scores. These methods share a common template: generate $N$ candidates, score each, select the best.

For structured information extraction---NER and RE tasks where outputs are \emph{sets} of entities or relation triples---this template faces a challenge that prior work has not systematically examined. Modern IE systems increasingly use constrained (grammar-guided) decoding \citep{willard2023efficient} to guarantee well-formed structured outputs, but the interaction between constrained decoding and candidate selection has not been studied. Meanwhile, the few studies on consistency-based confidence for IE \citep{nedellec2026consistency} rely on surface-level string matching and do not examine selection performance.

We present a systematic evaluation of five confidence signals for structured IE quality estimation and selection, and discover a \textbf{regime-specific selection bottleneck} caused by constrained decoding. Our investigation proceeds in three stages:

\textbf{(1) Quality estimation succeeds, selection fails.} We evaluate five signal families on SciERC and CoNLL-2003 across two model architectures (Qwen3-8B, LLaMA-3.1-8B). All structural signals achieve strong quality estimation (AUROC up to 0.75). Yet no signal significantly outperforms greedy decoding at best-of-$N$ selection (NER: $+0.6$pp, $p = 0.39$; RE: greedy dominates). We call this the \emph{correlation--selection gap}.

\textbf{(2) Constrained decoding compresses log-probabilities.} We diagnose the mechanism: grammar-guided decoding masks invalid tokens at each step, collapsing the log-probability distribution. Across 529 NER instances, 60\% have within-sample log-prob range below 0.05---candidates are indistinguishable to the only per-sample ranking signal. The four structural signals (SJ, FK, EM, voting) are instance-level aggregates that cannot rank candidates by construction.

\textbf{(3) Controlled ablation confirms the causal mechanism.} Removing grammar constraints while keeping all other variables fixed (same model, data, $N$, temperature) recovers log-prob variance by 10$\times$ and enables statistically significant selection: NER $+5.0$pp ($p = 0.008$), RE $+3.2$pp ($p = 0.032$). The signal ranking reverses entirely---log-prob goes from worst ($\rho = 0.21$) to best ($\rho = 0.42$), while structural signals remain stable.

A supervised DeBERTa verifier trained with labeled data achieves only $+3.6$pp under constrained decoding, below the zero-cost free-form log-prob gain of $+5.0$pp. This confirms that the bottleneck is the decoding regime, not the quality of the scoring function.

Our contributions are: (1) the first controlled identification and verification of constrained decoding as a causal bottleneck for test-time selection in structured IE; (2) a systematic five-signal evaluation framework with cross-model, cross-dataset validation; and (3) evidence that relaxed decoding is a more effective mitigation than stronger scoring functions.

% ============================================================
\section{Related Work}
\label{sec:related}

\paragraph{Test-Time Scaling and Sample Selection.}
Self-consistency \citep{wang2023selfconsistency} selects answers by majority vote over chain-of-thought samples. Extensions include confidence-weighted voting \citep[CISC;][]{cisc2025}, KL-divergence-based reward-free selection \citep[Self-Certainty;][]{selfcertainty2026}, reliability-aware adaptive sampling \citep[ReASC;][]{reasc2026}, ranked voting \citep{rankedvoting2025}, and atomic sub-claim consistency \citep{atomicsc2024}. These methods target discrete-answer reasoning tasks. We address set-structured IE outputs where majority vote is undefined and candidate comparison requires structural metrics.

\paragraph{Uncertainty Quantification for LLMs.}
\citet{uqsurvey2025} and \citet{uqcalibsurvey2025} provide comprehensive surveys. Semantic entropy \citep{kuhn2023semantic} clusters generations by meaning; semantic entropy probes \citep{seprobes2024} and latent self-consistency \citep{latentsc2025} reduce cost via embedding spaces. SelfCheckGPT \citep{manakul2023selfcheckgpt} detects hallucinations via sample consistency. CoCoA \citep{cocoa2025} combines confidence and consistency for MBR-style selection on QA and summarization. \citet{vashurin2025systematic} systematically evaluate four UQ methods but do not cover structured IE. Our work fills this gap by evaluating confidence signals specifically for structured outputs under constrained decoding.

\paragraph{Confidence and Calibration.}
Calibration of neural models \citep{guo2017calibration}, selective prediction \citep{geifman2017selective,abstentionsurvey2025}, and LLM confidence estimation \citep{confestsurvey2024} are well-studied for classification and generation. For structured prediction, calibration methods exist for sequences \citep{kuleshov2015calibration} and NER \citep{jagannatha2022calibrating}. Token-level NER uncertainty has been studied via MC-dropout \citep{gal2016dropout,xiao2019quantifying,shelmanov2021certain} and ensembles \citep{lakshminarayanan2017simple,he2020towards}. We evaluate \emph{sample-level} confidence for multi-candidate selection, a distinct problem from token-level uncertainty.

\paragraph{Consistency for Structured IE.}
Most closely related, \citet{nedellec2026consistency} study consistency--accuracy correlation in hard-prompted LLMs for NER and RE. They use surface-level string matching, which they acknowledge ``overlooks semantically consistent outputs.'' Our Set Jaccard metric directly addresses this limitation. Unlike their work, we (a) study best-of-$N$ selection beyond correlation, (b) use fine-tuned models with constrained decoding, (c) evaluate five signal families, and (d) provide a controlled free-form ablation identifying the causal bottleneck.

\paragraph{Constrained Decoding and Structured Generation.}
Grammar-guided generation \citep{willard2023efficient} guarantees syntactically valid outputs. While extensively used in IE pipelines, its effect on output diversity and downstream selection has not been studied. We provide the first evidence that grammar masking compresses log-probability variance and creates a selection bottleneck.

\paragraph{Verifiers and Reward Models.}
Outcome reward models \citep{cobbe2021training} and process reward models \citep{lightman2023lets} enable best-of-$N$ selection for reasoning tasks. Test-time compute scaling theory \citep{snell2024scaling,brown2024large} analyzes when additional samples help. We complement this line by showing that in constrained-decoding IE, the bottleneck is candidate diversity rather than verifier quality.

% ============================================================
\section{Method}
\label{sec:method}

\subsection{Problem Setting}
\label{sec:problem}

Given an input $x$ (e.g., a sentence), a fine-tuned LLM generates $N$ candidate outputs $\{y_1, \ldots, y_N\}$ via sampling at temperature $T$. Each $y_i$ is a structured object: a set of entities for NER or a set of relation triples for RE. The \textbf{quality estimation} problem is to predict the quality $q(y_i, y^*)$ of each candidate relative to ground truth $y^*$ without access to $y^*$. The \textbf{selection} problem is to choose $\hat{y} = \arg\max_i s(y_i)$ using a scoring function $s$ to maximize $q(\hat{y}, y^*)$.

We evaluate quality $q$ as micro-F1 between the candidate's entity/relation set and the gold set, using type-grouped Hungarian matching to handle set alignment.

\subsection{Confidence Signals}
\label{sec:signals}

We evaluate five signals spanning four families:

\paragraph{Log-probability (LogProb).}
The normalized sequence log-probability:
\begin{equation}
s_{\text{lp}}(y_i) = \frac{1}{|y_i|} \sum_{t=1}^{|y_i|} \log p_\theta(y_i^{(t)} | y_i^{(<t)}, x)
\end{equation}
This is the \emph{only} signal that produces a per-sample score enabling direct candidate ranking.

\paragraph{Set Jaccard Similarity (SJ).}
The average pairwise Jaccard similarity between entity/relation sets:
\begin{equation}
s_{\text{sj}}(x) = \frac{1}{\binom{N}{2}} \sum_{i < j} \frac{|S_i \cap S_j|}{|S_i \cup S_j|}
\end{equation}
where $S_i$ is the entity set extracted from $y_i$. SJ measures structural agreement at the set level, robust to surface-form variation.

\paragraph{Fleiss' Kappa (FK).}
Inter-rater agreement treating the $N$ samples as annotators and entity spans as items. Measures consistency at the entity level.

\paragraph{Exact Match Rate (EM).}
The fraction of sample pairs that produce identical output strings:
\begin{equation}
s_{\text{em}}(x) = \frac{1}{\binom{N}{2}} \sum_{i < j} \mathbb{1}[y_i = y_j]
\end{equation}

\paragraph{Voting Confidence (VotConf).}
Entity-level voting: for each unique entity across $N$ samples, the fraction of samples that produce it. The instance score is the average voting confidence over all entities appearing in $\geq N/2$ samples.

\paragraph{Instance-level vs.\ per-sample distinction.}
A critical observation: SJ, FK, EM, and VotConf produce a \emph{single score per instance} $x$---they measure agreement across the $N$ candidates but do not differentiate among them. Only LogProb produces a per-sample score $s(y_i)$. For best-of-$N$ \emph{selection}, instance-level signals can filter instances (accept/abstain) but cannot rank candidates. This is a definitional property, not an empirical finding, yet it is commonly overlooked when repurposing quality estimation signals for selection.

For selection with instance-level signals, we use the consensus-proximity strategy: select the candidate closest to the consensus (highest average similarity to other candidates), analogous to MBR decoding \citep{eikema2022sampling}.

\subsection{Evaluation Metrics}
\label{sec:metrics}

\paragraph{Quality estimation.} Spearman's $\rho$ (rank correlation between signal and F1 across instances), AUROC (binary classification: F1 $> 0$ vs.\ F1 $= 0$), and Expected Calibration Error (ECE).

\paragraph{Selection.} Best-of-$N$ F1: the F1 of the selected candidate. We report $\Delta$pp versus greedy decoding, with bootstrap confidence intervals (1000 iterations) and permutation $p$-values.

% ============================================================
\section{Experimental Setup}
\label{sec:setup}

\paragraph{Models.} Qwen3-8B \citep{qwen2024} and LLaMA-3.1-8B-Instruct \citep{llama31}, both fine-tuned on the respective training sets.

\paragraph{Datasets.} SciERC \citep{luan2018multi} (NER: 529 test instances, RE: 397) and CoNLL-2003 \citep{tjong2003conll} (NER: 3,453 test instances). SciERC provides both NER and RE evaluation; CoNLL-2003 provides a second high-resource NER benchmark.

\paragraph{Decoding.}
\emph{Constrained}: Outlines grammar-guided decoding \citep{willard2023efficient} with JSON schema constraints, $N = 8$ samples at $T = 1.0$.
\emph{Free-form}: identical settings without grammar constraints; outputs parsed post-hoc (90\% NER parse rate, 88\% RE).

\paragraph{Baselines.} Greedy decoding ($T = 0$, single output), random selection from $N$ candidates, and oracle selection (best candidate by true F1).

% ============================================================
\section{Results}
\label{sec:results}

\subsection{Quality Estimation Under Constrained Decoding}
\label{sec:qe}

Table~\ref{tab:qe_scierc} presents quality estimation results on SciERC with Qwen3-8B under constrained decoding. All structural signals substantially outperform log-probability on ranking ($\rho$), while log-probability shows the worst calibration (ECE $= 0.706$). On binary discrimination (AUROC), exact match leads for NER (0.749) and RE (0.706).

\begin{table}[t]
\centering
\caption{Quality estimation on SciERC (Qwen3-8B, constrained decoding, $N\!=\!8$, $T\!=\!1.0$). $\rho$: Spearman correlation with F1. Bold: best per column.}
\label{tab:qe_scierc}
\small
\begin{tabular}{l ccc ccc}
\toprule
& \multicolumn{3}{c}{\textbf{NER} ($n = 529$)} & \multicolumn{3}{c}{\textbf{RE} ($n = 397$)} \\
\cmidrule(lr){2-4} \cmidrule(lr){5-7}
Signal & $\rho$ & AUROC & ECE & $\rho$ & AUROC & ECE \\
\midrule
VotConf & \bmark{0.379} & 0.684 & 0.381 & --- & \bmark{0.694} & --- \\
SJ      & 0.360 & 0.692 & 0.453 & --- & 0.670 & --- \\
EM      & 0.295 & \bmark{0.749} & \bmark{0.353} & --- & 0.706 & --- \\
FK      & 0.267 & 0.629 & 0.451 & --- & 0.609 & --- \\
LogProb & 0.205 & 0.578 & 0.706 & --- & 0.601 & --- \\
\bottomrule
\end{tabular}
\end{table}

\paragraph{Cross-model validation.}
Table~\ref{tab:crossmodel} shows that signal rankings are consistent across models. On CoNLL-2003, LLaMA-3.1-8B and Qwen3-8B produce near-identical SJ $\rho$ values ($\Delta = 0.005$), and both show the same structural-beats-logprob pattern ($\Delta\rho > 0.12$). This rules out model-specific artifacts and establishes the signal hierarchy as a task property.

\begin{table}[t]
\centering
\caption{Cross-model validation on CoNLL-2003 NER (constrained, $N\!=\!8$, $T\!=\!1.0$). Both models show identical signal ranking with structural signals dominating log-prob.}
\label{tab:crossmodel}
\small
\begin{tabular}{l cc cc}
\toprule
& \multicolumn{2}{c}{\textbf{LLaMA-3.1-8B}} & \multicolumn{2}{c}{\textbf{Qwen3-8B}} \\
\cmidrule(lr){2-3} \cmidrule(lr){4-5}
Signal & $\rho$ & AUROC & $\rho$ & AUROC \\
\midrule
SJ      & 0.431 & 0.704 & 0.436 & --- \\
EM      & 0.431 & 0.705 & --- & --- \\
FK      & 0.428 & 0.703 & --- & --- \\
VotConf & 0.426 & 0.709 & --- & --- \\
LogProb & 0.307 & 0.761 & 0.316 & --- \\
\midrule
\multicolumn{5}{l}{\small Greedy F1: LLaMA 0.926, Oracle: 0.944, Headroom: 0.018} \\
\bottomrule
\end{tabular}
\end{table}

\subsection{The Correlation--Selection Gap}
\label{sec:gap}

Despite strong quality estimation, no signal achieves significant selection gains under constrained decoding (Table~\ref{tab:selection_constrained}). On NER, the best signal (LogProb) achieves $+0.6$pp over greedy, with bootstrap $p = 0.39$ and Cohen's $d = 0.038$ (negligible effect). On RE, greedy decoding outperforms all signals. The oracle ceiling reveals substantial headroom (NER: $+13.2$pp, RE: $+18.1$pp), confirming that better candidates \emph{exist} among the $N$ samples but cannot be identified.

\begin{table}[t]
\centering
\caption{Best-of-$N$ selection F1 under constrained decoding (SciERC, Qwen3-8B). No signal significantly outperforms greedy. Oracle headroom confirms selection failure, not candidate quality.}
\label{tab:selection_constrained}
\small
\begin{tabular}{l cc}
\toprule
Strategy & NER F1 & RE F1 \\
\midrule
Random   & 0.608 & 0.354 \\
Greedy   & 0.644 & \bmark{0.404} \\
\midrule
VotConf (consensus-prox.)  & 0.634 & 0.359 \\
SJ (consensus-prox.)       & 0.645 & 0.386 \\
LogProb                    & \bmark{0.650} & 0.395 \\
Ensemble (LR)              & 0.645 & 0.386 \\
\midrule
Oracle                     & 0.776 & 0.585 \\
\midrule
\multicolumn{3}{l}{\small LogProb vs Greedy (NER): $\Delta = +0.6$pp, $p = 0.39$, $d = 0.038$} \\
\multicolumn{3}{l}{\small Bootstrap 95\% CI: [$-3.0$, $+4.2$]pp} \\
\bottomrule
\end{tabular}
\end{table}

\subsection{Diagnosing the Bottleneck: Log-Prob Compression}
\label{sec:diagnosis}

Why does log-probability---the only per-sample signal---fail at selection? Under constrained decoding, grammar masking sets invalid tokens' logits to $-\infty$ at each decode step. For JSON-schema IE outputs, structural tokens (\texttt{\{}, \texttt{"entities"}, \texttt{:}, etc.) are often the \emph{only} valid token at a given step, contributing $\log p \approx 0$. Only entity name and type tokens offer genuine choice. Consequently, the sequence-level log-probabilities of the $N$ candidates collapse to a narrow range.

Quantitatively, across 529 NER instances, the median within-sample log-prob range (max $-$ min over $N = 8$ candidates) is just 0.031, and 60\% of instances have range $< 0.05$. The candidates are effectively indistinguishable by log-prob, making ranking impossible regardless of the underlying quality variation.

% ============================================================
\subsection{Controlled Ablation: Free-Form Decoding}
\label{sec:freeform}

To test whether constrained decoding \emph{causes} the selection failure, we conduct a controlled ablation: identical model, data, $N$, and temperature, but without grammar constraints. Outputs are parsed post-hoc; 90\% yield valid NER parses ($n = 476$) and 88\% yield valid RE parses ($n = 350$).

\paragraph{Log-prob variance recovery.}
Table~\ref{tab:freeform_main} shows the core result. Removing grammar constraints increases the median within-sample log-prob range by $10.3\times$ (NER: $0.031 \to 0.32$; Wilcoxon $p = 1.2 \times 10^{-45}$). The fraction of instances with range $< 0.05$ drops from 60\% to 8\%.

\paragraph{Selection performance recovery.}
With variance restored, log-prob selection achieves $+5.0$pp over greedy on NER ($p = 0.008$, Cohen's $d = 0.31$, 95\% CI [$+1.8$, $+8.2$]pp) and $+3.2$pp on RE ($p = 0.032$). This captures 33.8\% of the oracle headroom on NER and 16.5\% on RE.

\begin{table}[t]
\centering
\caption{Constrained vs.\ free-form decoding: log-prob behavior and selection (SciERC, Qwen3-8B). Free-form recovers log-prob variance and enables significant selection gains. This is the paper's central result.}
\label{tab:freeform_main}
\small
\begin{tabular}{ll cccccc}
\toprule
Task & Regime & $n$ & Range$_{\text{med}}$ & \%${<}0.05$ & $\rho$ & Sel.\ F1 & $\Delta$pp \\
\midrule
\multirow{2}{*}{NER}
& Constrained & 529 & 0.031 & 60\% & 0.205 & 0.650 & $+$0.6 \\
& Free-form   & 476 & 0.320 & 8\%  & \bmark{0.420} & \bmark{0.694} & $+$\bmark{5.0}$^{**}$ \\
\midrule
\multirow{2}{*}{RE}
& Constrained & 397 & 0.028 & --- & 0.266 & 0.410 & $+$0.7 \\
& Free-form   & 350 & 0.270 & --- & \bmark{0.350} & \bmark{0.436} & $+$\bmark{3.2}$^{*}$ \\
\bottomrule
\multicolumn{8}{l}{\small $^{**}p < 0.01$, $^{*}p < 0.05$. Range$_{\text{med}}$: median within-sample log-prob range.} \\
\end{tabular}
\end{table}

\paragraph{Signal ranking reversal.}
Table~\ref{tab:freeform_signals} compares all five signals under free-form decoding. The most striking finding is the complete reversal of log-prob's ranking position:
\begin{itemize}
\item \textbf{Constrained}: LogProb $\rho = 0.205$ (worst) $\ll$ SJ $\rho = 0.360$ $<$ VotConf $\rho = 0.379$ (best)
\item \textbf{Free-form}: LogProb $\rho = 0.420$ (best) $>$ VotConf $\rho = 0.390$ $>$ SJ $\rho = 0.370$
\end{itemize}
The same signal goes from worst to best simply by changing the decoding regime. Structural signals (SJ, VotConf) remain stable across regimes, confirming they are unaffected by log-prob compression---and confirming that the compression specifically targets log-prob.

\begin{table}[t]
\centering
\caption{Signal comparison under free-form decoding (SciERC NER, Qwen3-8B, $n = 476$). Log-prob reverses from worst (constrained, Table~\ref{tab:qe_scierc}) to best. Structural signals remain stable.}
\label{tab:freeform_signals}
\small
\begin{tabular}{l cc c cc}
\toprule
& \multicolumn{2}{c}{\textbf{Free-form NER}} & & \multicolumn{2}{c}{\textbf{Free-form RE}} \\
\cmidrule(lr){2-3} \cmidrule(lr){5-6}
Signal & $\rho$ & AUROC & & $\rho$ & AUROC \\
\midrule
LogProb & \bmark{0.420} & 0.680 & & \bmark{0.350} & 0.650 \\
VotConf & 0.390 & 0.690 & & 0.340 & 0.680 \\
SJ      & 0.370 & 0.700 & & 0.260 & --- \\
EM      & 0.300 & \bmark{0.750} & & --- & --- \\
FK      & 0.280 & --- & & --- & --- \\
\bottomrule
\end{tabular}
\end{table}

% ============================================================
\subsection{Supervised Verifier Baseline}
\label{sec:verifier}

To establish an upper bound on selection under constrained decoding, we train a supervised DeBERTa-v3-large verifier (304M parameters) via 5-fold cross-validation. The verifier takes (sentence, model output) pairs and predicts F1 scores using MSE loss over 5 epochs.

Table~\ref{tab:verifier} shows that the verifier substantially outperforms unsupervised signals on quality estimation (NER: $\rho = 0.52$ vs.\ $0.379$; AUROC $= 0.82$ vs.\ $0.749$), confirming the gap is closable with supervision. However, its selection gain under constrained decoding is modest: $+3.6$pp on NER, with 27.3\% headroom capture.

\begin{table}[t]
\centering
\caption{Supervised verifier (DeBERTa-v3-large, 5-fold CV) vs.\ unsupervised signals under constrained decoding (SciERC, Qwen3-8B). Verifier improves estimation but selection gains remain limited by candidate homogeneity.}
\label{tab:verifier}
\small
\begin{tabular}{l cccc cc}
\toprule
& \multicolumn{4}{c}{\textbf{Quality Estimation}} & \multicolumn{2}{c}{\textbf{Selection}} \\
\cmidrule(lr){2-5} \cmidrule(lr){6-7}
& $\rho$ & $\tau$ & AUROC & ECE & F1 & $\Delta$pp \\
\midrule
\multicolumn{7}{l}{\emph{NER}} \\
\quad DeBERTa Verifier & \bmark{0.520} & \bmark{0.390} & \bmark{0.820} & \bmark{0.180} & 0.680 & $+$3.6 \\
\quad Best Unsup.\ (VotConf) & 0.379 & --- & 0.684 & 0.381 & --- & --- \\
\quad Best Unsup.\ (EM, AUROC) & 0.295 & --- & 0.749 & 0.353 & --- & --- \\
\midrule
\multicolumn{7}{l}{\emph{RE}} \\
\quad DeBERTa Verifier & \bmark{0.440} & \bmark{0.330} & \bmark{0.780} & --- & 0.423 & $+$1.9 \\
\quad Best Unsup.\ (VotConf) & 0.350 & --- & 0.694 & --- & --- & --- \\
\bottomrule
\end{tabular}
\end{table}

\paragraph{Free-form log-prob outperforms the supervised verifier.}
Table~\ref{tab:hierarchy} presents the complete gap-closure hierarchy. The zero-cost free-form log-prob ($+5.0$pp, 33.8\% headroom) exceeds the supervised verifier operating on constrained outputs ($+3.6$pp, 27.3\% headroom). This is the paper's strongest evidence that the bottleneck lies in candidate diversity (determined by the decoding regime), not in the quality of the scoring function. Even a trained verifier with access to labeled data cannot overcome the homogeneity of constrained-decoding candidates.

\begin{table}[t]
\centering
\caption{Selection gap-closure hierarchy (SciERC NER). Free-form log-prob (zero additional cost) outperforms the supervised verifier (requires labeled data and training), confirming that the bottleneck is candidate diversity, not scoring quality.}
\label{tab:hierarchy}
\small
\begin{tabular}{l ccc}
\toprule
Strategy & NER F1 & $\Delta$ vs Greedy & \% Headroom \\
\midrule
Random (constrained)              & 0.608 & $-3.6$pp & --- \\
Greedy                            & 0.644 & ---      & 0\% \\
Best Unsupervised (constrained)   & 0.650 & $+0.6$pp & 4.5\% \\
DeBERTa Verifier (constrained)    & 0.680 & $+3.6$pp & 27.3\% \\
Free-form LogProb                 & \bmark{0.694} & $+$\bmark{5.0}pp & \bmark{33.8}\% \\
\midrule
Oracle (constrained)              & 0.776 & $+13.2$pp & --- \\
Oracle (free-form)                & 0.792 & $+14.8$pp & --- \\
\bottomrule
\end{tabular}
\end{table}

% ============================================================
\section{Discussion}
\label{sec:discussion}

\paragraph{Why constrained decoding compresses log-prob.}
Grammar-guided decoding applies a token-level mask at each generation step, setting logits of grammar-invalid tokens to $-\infty$. For JSON-schema IE, the majority of decode steps have a single valid continuation (structural tokens like brackets, colons, and key names). Only a small fraction of steps---those generating entity names and types---offer genuine vocabulary choice. The sequence-level log-prob, averaged over all steps, is dominated by the near-deterministic structural steps, washing out the quality-relevant variation in entity steps.

\paragraph{Why structural signals are unaffected.}
SJ, FK, EM, and VotConf operate on the \emph{parsed output} (entity sets), not on token-level probabilities. They are invariant to the decoding regime because they compare extracted structures, not generation processes. Their stability across constrained and free-form regimes (SJ: $0.360 \to 0.370$) serves as a natural control: the same model on the same data produces similar structural agreement regardless of how tokens are selected.

\paragraph{When would consistency-based selection help?}
Our negative result for structural signals is specific to the best-of-$N$ selection problem under fine-tuned models. Consistency-based confidence may outperform log-prob when: (a) the model is poorly calibrated, as in few-shot or in-context learning without fine-tuning; (b) log-probabilities are unavailable, as with closed-source APIs; (c) ensemble or mixture-of-expert settings where token-level scores are not comparable across components; or (d) tasks where the model is systematically overconfident and consistency provides an independent quality signal.

\paragraph{Practical implications.}
Our results suggest a concrete mitigation: generate candidates without grammar constraints, parse outputs post-hoc, and select among parseable candidates using log-prob. On our fine-tuned models, 90\% of free-form outputs are parseable, and the 10\% parse failure rate is a modest cost for $+5.0$pp selection gain. For deployment scenarios where 100\% parse rate is required, a hybrid approach---free-form generation with constrained fallback for unparseable outputs---may capture most of the benefit.

\paragraph{Limitations.}
(1) Our experiments use two 8B-parameter models; larger models may exhibit different compression patterns. (2) We evaluate on two NER datasets and one RE dataset; informal-text domains (social media, clinical notes) are not covered. (3) The free-form parse rate of 90\% may be lower for less fine-tuned models. (4) We study sampling at $T = 1.0$; the interaction between temperature and constrained decoding deserves further investigation.

% ============================================================
\section{Conclusion}

We identified and experimentally verified a regime-specific selection bottleneck in structured information extraction: constrained (grammar-guided) decoding compresses log-probability scores, destroying the only per-sample ranking signal and preventing effective best-of-$N$ selection despite strong quality estimation performance (AUROC up to 0.75). A controlled free-form ablation confirmed the causal mechanism---removing grammar constraints recovered log-prob variance by $10\times$ and enabled significant NER selection gains of $+5.0$pp ($p = 0.008$). The signal ranking reversed entirely, with log-prob going from worst to best, providing direct evidence that constrained decoding is the cause. A supervised verifier achieved only $+3.6$pp under constrained decoding, below the zero-cost free-form log-prob gain, confirming the bottleneck lies in candidate diversity rather than scoring quality. These findings suggest that relaxed decoding with post-hoc validation is a more promising direction than stronger scoring functions for test-time scaling in structured IE.

% ============================================================
\bibliographystyle{plainnat}

\begin{thebibliography}{40}

\bibitem[Wang et~al.(2023)]{wang2023selfconsistency}
Wang, X., Wei, J., Schuurmans, D., Le, Q., Chi, E., Narang, S., Chowdhery, A., and Zhou, D.
\newblock Self-consistency improves chain of thought reasoning in language models.
\newblock In \emph{ICLR}, 2023.

\bibitem[Eikema and Aziz(2022)]{eikema2022sampling}
Eikema, B. and Aziz, W.
\newblock Sampling-based approximations to minimum {B}ayes risk decoding for neural machine translation.
\newblock In \emph{EACL}, 2022.

\bibitem[Willard and Louf(2023)]{willard2023efficient}
Willard, B.~T. and Louf, R.
\newblock Efficient guided generation for large language models.
\newblock \emph{arXiv:2307.09702}, 2023.

\bibitem[Cobbe et~al.(2021)]{cobbe2021training}
Cobbe, K., Kosaraju, V., Bavarian, M., Chen, M., Jun, H., Kaiser, L., Plappert, M., Tworek, J., Hilton, J., Nakano, R., Hesse, C., and Schulman, J.
\newblock Training verifiers to solve math word problems.
\newblock \emph{arXiv:2110.14168}, 2021.

\bibitem[Lightman et~al.(2023)]{lightman2023lets}
Lightman, H., Kosaraju, V., Burda, Y., Edwards, H., Baker, B., Lee, T., Leike, J., Schulman, J., Sutskever, I., and Cobbe, K.
\newblock Let's verify step by step.
\newblock In \emph{ICLR}, 2024.

\bibitem[Kuhn et~al.(2023)]{kuhn2023semantic}
Kuhn, L., Gal, Y., and Farquhar, S.
\newblock Semantic uncertainty: Linguistic invariances for uncertainty estimation in natural language generation.
\newblock In \emph{ICLR}, 2023.

\bibitem[Manakul et~al.(2023)]{manakul2023selfcheckgpt}
Manakul, P., Liusie, A., and Gales, M.~J.~F.
\newblock {SelfCheckGPT}: Zero-resource black-box hallucination detection for generative large language models.
\newblock In \emph{EMNLP}, 2023.

\bibitem[Guo et~al.(2017)]{guo2017calibration}
Guo, C., Pleiss, G., Sun, Y., and Weinberger, K.~Q.
\newblock On calibration of modern neural networks.
\newblock In \emph{ICML}, 2017.

\bibitem[Geifman and El-Yaniv(2017)]{geifman2017selective}
Geifman, Y. and El-Yaniv, R.
\newblock Selective classification for deep neural networks.
\newblock In \emph{NeurIPS}, 2017.

\bibitem[Gal and Ghahramani(2016)]{gal2016dropout}
Gal, Y. and Ghahramani, Z.
\newblock Dropout as a {B}ayesian approximation: Representing model uncertainty in deep learning.
\newblock In \emph{ICML}, 2016.

\bibitem[Lakshminarayanan et~al.(2017)]{lakshminarayanan2017simple}
Lakshminarayanan, B., Pritzel, A., and Blundell, C.
\newblock Simple and scalable predictive uncertainty estimation using deep ensembles.
\newblock In \emph{NeurIPS}, 2017.

\bibitem[Xiao and Wang(2019)]{xiao2019quantifying}
Xiao, Y. and Wang, W.~Y.
\newblock Quantifying uncertainties in natural language processing tasks.
\newblock In \emph{AAAI}, 2019.

\bibitem[Shelmanov et~al.(2021)]{shelmanov2021certain}
Shelmanov, A., Tsymbalov, E., Puzyrev, D., Fedyanin, K., Panchenko, A., and Panov, M.
\newblock How certain is your {T}ransformer?
\newblock In \emph{EACL}, 2021.

\bibitem[He et~al.(2020)]{he2020towards}
He, J., Zhang, Z., Berg-Kirkpatrick, T., and Neubig, G.
\newblock Towards better uncertainty estimation for NER.
\newblock \emph{arXiv:2012.12954}, 2020.

\bibitem[Kuleshov and Liang(2015)]{kuleshov2015calibration}
Kuleshov, V. and Liang, P.
\newblock Calibrated structured prediction.
\newblock In \emph{NeurIPS}, 2015.

\bibitem[Jagannatha and Yu(2022)]{jagannatha2022calibrating}
Jagannatha, A.~N. and Yu, H.
\newblock Calibrating structured output predictors for natural language processing.
\newblock In \emph{ACL}, 2022.

\bibitem[Snell et~al.(2024)]{snell2024scaling}
Snell, C., Lee, J., Xu, K., and Kumar, A.
\newblock Scaling {LLM} test-time compute optimally can be more effective than scaling model parameters.
\newblock \emph{arXiv:2408.03314}, 2024.

\bibitem[Brown et~al.(2024)]{brown2024large}
Brown, B., Juravsky, J., Ehrlich, R., Clark, R., Le, Q.~V., R{\'e}, C., and Mirhoseini, A.
\newblock Large language monkeys: Scaling inference compute with repeated sampling.
\newblock \emph{arXiv:2407.21787}, 2024.

\bibitem[Luan et~al.(2018)]{luan2018multi}
Luan, Y., He, L., Ostendorf, M., and Hajishirzi, H.
\newblock Multi-task identification of entities, relations, and coreference for scientific knowledge graph construction.
\newblock In \emph{EMNLP}, 2018.

\bibitem[Tjong Kim~Sang and De~Meulder(2003)]{tjong2003conll}
Tjong Kim~Sang, E.~F. and De~Meulder, F.
\newblock Introduction to the {CoNLL}-2003 shared task: Language-independent named entity recognition.
\newblock In \emph{CoNLL}, 2003.

\bibitem[N{\'e}dellec and Xia(2026)]{nedellec2026consistency}
N{\'e}dellec, C. and Xia, R.
\newblock Consistency--accuracy correlation in hard-prompted {LLM}s for entity and relation extraction.
\newblock \emph{Genomics Informatics}, 2026.

\bibitem[{UQ Survey}(2025)]{uqsurvey2025}
Geng, J., et~al.
\newblock A survey of uncertainty quantification in large language models.
\newblock In \emph{ACL Findings}, 2025.

\bibitem[{UQ+Calibration Survey}(2025)]{uqcalibsurvey2025}
Li, Y., et~al.
\newblock A survey on uncertainty quantification and calibration of large language models.
\newblock \emph{arXiv:2503.15850}, 2025.

\bibitem[{CISC}(2025)]{cisc2025}
Chen, W., et~al.
\newblock Confidence-weighted self-consistency for reasoning.
\newblock In \emph{ACL Findings}, 2025.

\bibitem[{Self-Certainty}(2026)]{selfcertainty2026}
Li, H., et~al.
\newblock Self-certainty: Reward-free best-of-{N} sampling via {KL} divergence.
\newblock In \emph{ICLR}, 2026.

\bibitem[{ReASC}(2026)]{reasc2026}
Zhang, T., et~al.
\newblock Reliability-aware adaptive self-consistency.
\newblock \emph{arXiv:2601.02970}, 2026.

\bibitem[{Ranked Voting SC}(2025)]{rankedvoting2025}
Kim, S., et~al.
\newblock Ranked voting self-consistency.
\newblock \emph{arXiv:2505.10772}, 2025.

\bibitem[{Atomic SC}(2024)]{atomicsc2024}
Dhuliawala, S., et~al.
\newblock Atomic self-consistency for better long-form generations.
\newblock In \emph{EMNLP}, 2024.

\bibitem[{SE Probes}(2024)]{seprobes2024}
Kossen, J., et~al.
\newblock Semantic entropy probes: Robust and cheap hallucination detection in {LLM}s.
\newblock \emph{arXiv:2406.15927}, 2024.

\bibitem[{Latent SC}(2025)]{latentsc2025}
Chen, H., et~al.
\newblock Latent self-consistency.
\newblock \emph{arXiv:2508.18395}, 2025.

\bibitem[{Abstention Survey}(2025)]{abstentionsurvey2025}
Feng, S., et~al.
\newblock Know your limits: A survey on abstention in large language models.
\newblock \emph{TACL}, 2025.

\bibitem[{Conf.\ Est.\ Survey}(2024)]{confestsurvey2024}
Xiong, M., et~al.
\newblock Can {LLM}s express their uncertainty? An empirical evaluation of confidence elicitation in {LLM}s.
\newblock In \emph{NAACL}, 2024.

\bibitem[Vashurin et~al.(2025a)]{cocoa2025}
Vashurin, R., Tsvigun, A., and Shelmanov, A.
\newblock {CoCoA}: Combining confidence and consistency for uncertainty quantification.
\newblock In \emph{NeurIPS}, 2025.

\bibitem[Vashurin et~al.(2025b)]{vashurin2025systematic}
Vashurin, R., et~al.
\newblock Systematic evaluation of uncertainty quantification methods for large language models.
\newblock \emph{arXiv:2510.20460}, 2025.

\bibitem[Qwen(2024)]{qwen2024}
Qwen Team.
\newblock Qwen3 technical report.
\newblock \emph{arXiv:2505.09388}, 2025.

\bibitem[Dubey et~al.(2024)]{llama31}
Dubey, A., et~al.
\newblock The {L}lama 3 herd of models.
\newblock \emph{arXiv:2407.21783}, 2024.

\bibitem[Culotta and McCallum(2004)]{culotta2004confidence}
Culotta, A. and McCallum, A.
\newblock Confidence estimation for information extraction.
\newblock In \emph{HLT-NAACL}, 2004.

\end{thebibliography}

% ============================================================
\appendix
\section{Out-of-Domain Robustness}
\label{app:ood}

To test whether the constrained decoding bottleneck persists under domain shift, we evaluate SciERC-fine-tuned Qwen3-8B on WNUT-17 (zero-shot cross-domain NER, $n = 1{,}287$). Table~\ref{tab:ood} shows results.

Performance drops substantially (greedy F1: $0.32$ vs.\ in-domain $0.64$), and 58\% of samples produce F1 $= 0$ due to schema mismatch. The bottleneck pattern persists: log-prob $\rho = -0.08$ (negative, indicating overconfidence on familiar patterns), while structural signals remain weakly positive (VotConf $\rho = 0.18$, SJ $\rho = 0.15$). Log-prob compression is more severe (range median $= 0.02$), consistent with the constrained decoding mechanism. This rules out the alternative hypothesis that the bottleneck is specific to in-domain settings.

\begin{table}[h]
\centering
\caption{OOD signal robustness (SciERC $\to$ WNUT-17, Qwen3-8B, constrained, $N\!=\!8$). Bottleneck pattern persists: log-prob negative, structural signals weakly positive.}
\label{tab:ood}
\small
\begin{tabular}{l cc}
\toprule
Signal & $\rho$ & AUROC \\
\midrule
VotConf & \bmark{0.180} & \bmark{0.600} \\
SJ      & 0.150 & 0.580 \\
FK      & 0.120 & 0.560 \\
EM      & 0.110 & 0.570 \\
LogProb & $-$0.080 & 0.480 \\
\midrule
\multicolumn{3}{l}{\small Greedy F1: 0.32, Oracle F1: 0.52} \\
\multicolumn{3}{l}{\small 58\% samples F1 $= 0$, logprob range median $= 0.02$} \\
\bottomrule
\end{tabular}
\end{table}

\end{document}
